\section{Introduction}
\label{sec:intro}

Behavioral biometrics — the study of \emph{how} individuals interact with
digital devices — offers a passive, non-invasive window into cognitive and
motor state~\cite{monrose2000keystroke}.
Among the richest behavioral signals is \emph{keystroke dynamics}: the
millisecond-level rhythm, speed, error patterns, and per-key latency produced
during typing.
These signals capture aspects of neuromuscular coordination that may reflect
fatigue, stress, or longer-term behavioral change, making them attractive for
longitudinal personal monitoring.

Existing keystroke dynamics systems predominantly focus on
\emph{authentication}, comparing a user's current pattern against a stored
population template to verify identity~\cite{killourhy2009comparing}.
This framing treats the typing template as static and optimizes for
\emph{inter-person} discrimination.
A complementary and less-studied problem is detecting \emph{within-person
behavioral drift} over time: has this individual's typing changed relative to
their own established baseline?
This question is directly relevant to cognitive health monitoring, workplace
fatigue detection, and human-computer interaction research, but standard
population-norm approaches fail at it by design — a consistently slow typist
would be permanently flagged as anomalous.

\paragraph{The N-of-1 Paradigm.}
The N-of-1 clinical trial design~\cite{lillie2011nof1} treats each
individual as their own control.
Applying this philosophy to deep learning, we train a model \emph{exclusively}
on a single user's longitudinal data, learning what regularity looks like
\emph{for that person}, and flagging deviations from their own history.
This avoids all population-norm assumptions and is inherently personalized.

\paragraph{NeuroLink 3.0.}
We present NeuroLink 3.0, a system that learns a compact 10-dimensional
\emph{Muscle Memory Signature} from 42,501 personal typing lessons using an
undercomplete autoencoder.
High reconstruction error on a new lesson signals behavioral deviation.
This per-lesson error, aggregated into a 7-day rolling Longitudinal Drift
Index (LDI), forms the core anomaly signal.
Differential privacy via DP-SGD~\cite{abadi2016deep} ensures the model cannot
memorize individual lessons.
\Cref{fig:dashboard} shows the full system output over 435 deployment days.

\paragraph{Contributions.} This paper makes the following contributions:

\begin{enumerate}
  \item \textbf{N-of-1 undercomplete autoencoder} trained on 17 months of
        personal keystroke data with a 10-dimensional bottleneck that learns
        individual behavioral regularity without reference to population norms.

  \item \textbf{Residual feature engineering} for EV and ESRP that resolves a
        structural reconstruction ceiling caused by using rolling-mean
        formulations — features that cannot be reconstructed from a single-lesson
        snapshot at inference time.

  \item \textbf{Differential privacy} via Opacus DP-SGD~\cite{yousefpour2021opacus}
        with RDP accounting, achieving \(\varepsilon{=}7.9983\) at
        \(\delta{=}10^{-5}\) — providing a formal privacy guarantee that the
        behavioral baseline cannot leak raw lesson data.

  \item \textbf{Bootstrap confidence interval alerting}: 2,000 bootstrap
        resamples of training LDI values produce a 95\% CI band enabling
        three-tier classification (HIGH\,/\,MARGINAL\,/\,NORMAL), reducing
        false positives compared to a fixed point threshold.

  \item \textbf{Empirical Ghost V3}: a drift-adjusted typing pace predictor
        whose scale factor is derived from a linear regression of daily LDI
        ratio against observed WPM across the full 42,501-lesson dataset,
        replacing a hand-tuned constant.
\end{enumerate}
